\mode*
\section{Theoretical background}
\label{background}

The background necessary to understand the rest of this paper is
phenomenography, variation theory, and the approaches to learning.
We briefly recap these concepts; for a fuller treatment in the context of
introductory programming, see the companion paper
\autocite{VTProgMisconceptions}.

\subsection{Phenomenography}

\begin{frame}
  \textcquote[pp.~112--113]{NCOL}{%
    \textins{U}nderstanding does not cause acts; instead, acts express
    understanding%
  }.
\end{frame}
This quote captures the core of phenomenography and the essence of the type of
data we need: we cannot observe understanding directly, but we can observe
acts---and debugging is very much a sequence of acts.

Phenomenography rests on two premises: that reality is always experienced
rather than apprehended directly, and that people interpret the significant
aspects of that reality in different ways. From this follows a distinction
between two perspectives. A \emph{first-order} perspective describes various
aspects of the world as they are taken to be; a \emph{second-order} perspective
instead describes how people \emph{experience} those aspects. Phenomenography
adopts the second-order perspective: its object of study is not the world itself
but the qualitatively different ways in which people experience it
\autocite{Phenomenography}.

\mode<presentation>{%
\begin{frame}
  \begin{block}<1-2>{Phenomenography~\autocite{Phenomenography}}
    \begin{itemize}
      \item Reality is experienced.
      \item People interpret significant aspects of reality.
    \end{itemize}
    \pause
    \begin{description}
      \item[First-order perspective] describes various aspects of the world.
      \item[Second-order perspective] describes people's experiences of various
        aspects of the world---phenomenography.
    \end{description}
  \end{block}

  \pause

  \begin{remark}
    \textcquote[pp.~112--113]{NCOL}{%
      \textins{U}nderstanding does not cause acts; instead, acts express
      understanding%
    }.
  \end{remark}
\end{frame}%
}

For debugging, the second-order perspective is what matters: the bug exists
precisely in the gap between the program as it is (first order) and the
program as the programmer experiences it (second order).

\subsection{Variation theory}

\begin{frame}
  \blockcquote[p.~44]{NCOL}{%
    Imagine that you live in an entirely green world!
    Everything is equally green and is always green.
    Exactly the same nuance, no difference at all.
    Do you think that you would notice the greenness of the grass in that
    world, and would you know what \enquote{green} is at all?
    \textelp{}
    Would it be possible for you to learn to see greenness if a teacher coming
    from another world tried to help you by pointing to all kind of green
    things \textelp{} saying \enquote{green} each time?%
  }
\end{frame}

These ideas form a chain: there is no learning without discernment, and no
discernment without variation. Learning therefore depends on discernment, which
in turn depends on the learner experiencing a pattern of variation
\autocites[Ch.~3]{NCOL}{VariationTheory}.

\mode<presentation>{%
\begin{frame}
  \begin{block}{Variation theory~\autocite{VariationTheory}}
    \vspace{-0.5em}
    \[
      \text{learning}
      \quad\leftarrow\quad
      \text{discernment}
      \quad\leftarrow\quad
      \text{variation}
    \]
  \end{block}

  \pause

  \begin{remark}
    \begin{itemize}
      \item There must be a pattern of variation to experience.
      \item \emph{This pattern must be experienced} for learning to happen.
    \end{itemize}
  \end{remark}
\end{frame}%
}

Each educational objective can be divided into \emph{aspects}
\autocite[Ch.~2]{NCOL}.
Aspects that the learner has not yet discerned are the \emph{critical}
aspects.
To discern a critical aspect the learner must experience variation in the
dimension of that aspect.
The patterns of variation are (with the example of learning what the colour
green is):

\begin{description}
  \item[Contrast] Vary the critical aspect, keep the rest invariant: show a
    green circle next to a blue circle.
    This opens up the dimension of variation (colour) and separates the
    critical feature (green) from other features.
  \item[Generalisation] Vary the non-critical aspects, keep the critical
    aspect invariant: show several green figures of different shapes.
    This separates the critical aspect from the incidental ones.
    (Called induction if it comes before contrast.)
  \item[Fusion] Vary the aspects in focus together: show figures varying in
    both colour and shape at once, bringing together what contrast and
    generalisation separated.
\end{description}

\begin{frame}\label<1>{vtsummary}
  \begin{figure}
    \begin{subfigure}{0.3\columnwidth}
      \centering
      \includegraphics{figs/contrast-color.tikz}
      \caption{Contrast}
    \end{subfigure}
    \hfill
    \begin{subfigure}{0.3\columnwidth}
      \centering
      \includegraphics{figs/generalization-color.tikz}
      \caption{Generalization}
    \end{subfigure}
    \hfill
    \begin{subfigure}{0.3\columnwidth}
      \centering
      \includegraphics{figs/fusion-color.tikz}
      \caption{Fusion}
    \end{subfigure}
    \caption{%
      Illustrating the patterns of variation for aspects color and shape.
    }\label{fig:vt-patterns}
  \end{figure}

  \begin{onlyenv}<1>
    \mode<presentation>{%
      \begin{example}
        \begin{itemize}
          \item Critical aspect: colour.
          \item Critical feature: green; contrasted against blue.
          \item Non-critical aspects: shape, order.
          \item Non-critical features: circle, square; order of pairs.
        \end{itemize}
      \end{example}%
    }
  \end{onlyenv}
\end{frame}

In terms of \cref{fig:vt-patterns}: the critical aspect is colour, and the critical
feature---the value the learner is to discern---is green, which the
contrast sets against blue.
Shape and the order of the pairs are non-critical aspects, whose features
(circle and square, and the order in which the pairs appear) may vary freely.
The generalisation pattern exploits exactly this freedom: the shape varies
while green is kept invariant.
In the fusion pattern, finally, colour and shape vary together---a green
square against a yellow circle---so the learner meets the critical feature
amid simultaneous variation in both dimensions.

One can vary the order of these steps, but the contrast should come first,
and \citeauthor{NCOL} motivates why \autocite[pp.~45--47]{NCOL}.
Generalisation before contrast---induction---asks the learner to notice
what several green things have in common; but that cannot open a dimension
the learner has never discerned:
\textcquote[p.~46]{NCOL}{if you cannot see the greenness of any of three
green, though otherwise different, things, you cannot see what they have in
common, either}.
Of two teachers who juxtapose sets of objects to make greenness discernible,
\textcquote[p.~47]{NCOL}{\textins{o}ne is betting on induction, the other on
contrast. The latter wins}.
There is also a cognitive-load motivation hiding in the same passage:
discernment comes from
\textcquote[p.~46]{NCOL}{the simultaneous experience of things that differ in
color but are the same otherwise}---and instances that are not juxtaposed
side by side must be juxtaposed in the learner's mind.
Starting with the contrast puts the differing values next to each other;
any other order makes the learner carry the earlier instances in memory
until the difference can be experienced.

\subsection{Approaches to learning}
\label{approaches}

The distinction between \emph{deep} and \emph{surface} approaches to learning
originates with \textcite{MartonSaljo1976I}, who studied university students
reading prose texts.
Students who engaged in what was termed deep-level processing directed their
attention towards what the text was \emph{about}---its meaning---whereas
students engaging in surface-level processing directed their attention towards
the text itself, towards being able to reproduce it.
The two groups reached qualitatively different learning outcomes: the study
describes differences in \emph{what} is learned, rather than in how much.
The approach is not a fixed trait of the student:
\textcite{MartonSaljo1976II} showed that students adapt their way of learning
to their conception of what the task requires of them.

\mode<presentation>{%
\begin{frame}
  \begin{description}
    \item[Deep approach] attention on the \emph{meaning}---what the text is
      about.
    \item[Surface approach] attention on the \emph{sign}---the text itself,
      reproducing it.
  \end{description}

  \pause

  \begin{remark}
    \begin{itemize}
      \item Qualitatively different outcomes: \emph{what} is learned, not how
        much.
      \item Not a fixed trait: students adapt the approach to the perceived
        task.
    \end{itemize}
  \end{remark}
\end{frame}%
}

The approaches can be measured: the revised two-factor Study Process
Questionnaire (R-SPQ-2F) assesses deep and surface approaches with ten items
per approach and is designed to be usable by teachers with their own students
\autocite{RSPQ2F}.
We return to this instrument in \cref{method}.

In \citetitle{NCOL}, \citeauthor{NCOL} reinterprets the approaches to learning
in terms of variation theory \autocite[pp.~142--146]{NCOL}.
Already in the original studies the two approaches differed in terms of
variation and invariance: those adopting a deep approach
\textcquote[p.~142]{NCOL}{tried to find the main point in the text, and the
words for capturing it, by trying alternatives}, whereas those adopting a
surface approach tried to make sure that they would remember what they had
read.
Marton concludes that to develop a powerful way of seeing something, the
learner must decompose the object of learning and bring it together again,
and towards that end
\textcquote[p.~145]{NCOL}{the learner has to create the necessary patterns of
variation and invariance. This is the deep approach to learning in terms of
the theory of variation}.
He summarises: \textcquote[p.~151]{NCOL}{the more powerful, deep approach to
learning is more varied}, and the surface approach less varied.
This reinterpretation is the pivot of our argument in
\cref{debugging-as-learning}.

\subsection{The schooling the students bring with them}
\label{schooling}

The approach a student takes is not settled by the course alone: students
adapt it to what they take the task to require
\autocite{MartonSaljo1976II}, and by the time they arrive they have spent
twelve years being told what studying requires.
For a Swedish cohort, one feature of those twelve years is recorded
nationally and needs no self-report---the type of principal
(\emph{huvudman}) that ran each school unit, whether the municipality
(\emph{kommunal}) or an independent principal (\emph{enskild}, a
\emph{friskola}).
There is a Swedish literature on what that type is associated with, and it
points somewhere interesting.

Independent schools grade more generously than municipal schools relative
to externally measured knowledge.
\Textcite{HinnerichVlachos2017} had the national tests of upper-secondary
students re-graded by external assessors and compared the two gradings of
the same tests: the internal grading of voucher schools was about
\(0.133\) standard deviations more generous
\autocite[Table~2]{HinnerichVlachos2017}.
Their students \textcquote[p.~40]{HinnerichVlachos2017}{come out around
0.15 standard deviations ahead} on the internal evaluations while scoring
no better on the external ones, and the authors conclude that they do so
\textcquote[p.~42]{HinnerichVlachos2017}{despite not having built more or
better human capital}.
Skolverket's analyses of grades against the national tests find the same
in year~9 \autocite[p.~37]{Skolverket2019Betyg} and in the gymnasium,
where about half of the independent schools' grade advantage would
disappear if the two sectors graded alike
\autocite[p.~42]{Skolverket2020Betyg}.
\Textcite{EdmarkPersson2021} find that independent schools set a course
grade above the national test result more often than municipal schools do,
and below it no more often, and that the effect vanishes in
mathematics---the test that leaves the least room for discretion.
\Textcite{WikstromWikstrom2005} reported the same pattern against an
external test two decades earlier.
The evidence is strongest at the upper-secondary level; at the compulsory
level \textcite{BohlmarkLindahl2015} find no differential grade inflation
at all.

This is a claim about grades, not about knowledge, and the distinction has
to be kept.
In the same re-grading study the difference in the externally graded
results themselves is small, negative and not statistically significant
\autocite[Table~2]{HinnerichVlachos2017}, and
\textcite{BohlmarkLindahl2015} find that the voucher reform improved
average outcomes, mainly through competition between schools.
The literature does not say that students from independent schools know
less; it says that their grades run ahead of what external assessment
shows.

The discrepancy appears again at the transition to higher education.
Skolverket's register study of that transition reports that students who
graduated from an independent gymnasium complete fewer of their first-year
higher-education credits than students from municipal schools---61 against
66 per cent complete at least three quarters of the credits they
registered for, and 21 against 16 per cent complete less than half
\autocite[p.~30]{Skolverket2018Hogskola}---although they arrive with
higher grades, and although the difference survives controls for sex,
background, parental education, upper-secondary grades and choice of study
(odds ratios between \(1.58\) and \(2.07\) for the risk of completing less
than half, all significant at the \(0.1\) per cent level)
\autocite[p.~39]{Skolverket2018Hogskola}.
A joint successor study replicates this across eight cohorts and rules out
the mechanical explanation that these students merely register for more
credits \autocite[pp.~36--37]{UKASkolverket2024}.
Both are agency register studies rather than research findings, and
neither identifies a cause; the later one says outright that its
statistics cannot explain what the differences are due to
\autocite[p.~47]{UKASkolverket2024}.
At the other margin the sign reverses:
\textcite{EdmarkPersson2021} find students from independent schools more
likely to reach higher education in the first place.
They get there more often, and keep up less well once there.

Why they keep up less well is not answered anywhere.
The literature offers no mechanism connecting the huvudman type to how
students go about their learning; we searched for one twice and found
nothing (\cref{app:search}).
What the grading literature does offer is an explanation on the supply
side: \textcite[pp.~40--41]{HinnerichVlachos2017} read their own results
as schools responding to differences in educational preferences, some of
them not focusing on core subject knowledge---a segmented market rather
than a difference in how the students learn.
Our conjecture is a different one, and it is ours rather than the
literature's.
A grade that runs ahead of demonstrated knowledge is what one would expect
where reproducing the assessed material suffices, and attending to the
material in order to reproduce it is precisely the surface approach
(\cref{approaches}).
If the assessment a student met for twelve years rewarded reproduction,
the approach they bring to university may be the one that reproduction
rewarded.
That is a conjecture about students rather than about schools, and this
study is in a position to look at it: we measure the approach to learning
anyway, and the schools the students came from can be looked up.
This is what \cref{rq:school} asks.

\mode<presentation>{%
\begin{frame}
  \begin{block}{Swedish schools: grades and knowledge}
    \begin{itemize}
      \item Independent schools (\emph{friskolor}) grade more generously
        against external assessment \autocite{HinnerichVlachos2017}.
      \item Not that their students know less---the externally graded
        results do not differ significantly
        \autocite{HinnerichVlachos2017}.
      \item Their graduates complete fewer first-year university credits
        than their grades predict
        \autocite{Skolverket2018Hogskola,UKASkolverket2024}.
    \end{itemize}
  \end{block}

  \pause

  \begin{idea}
    \begin{itemize}
      \item Nobody has asked whether the approach to learning is the
        mechanism.
      \item Our conjecture: reproduction rewarded \(\to\) surface approach
        brought along.
      \item The rival: schools segmenting a market, not students learning
        differently.
    \end{itemize}
  \end{idea}
\end{frame}%
}
